\section{面向组合查询的域内任务数据}
\label{sec:domain-task-data}

\subsection{从参数读法到可执行答案}

专业资料中的选型与配置问题通常需要同时处理多个参数。例如，判断某一环境温度是否符合设备说明，需要区分扫描器与显示部件；筛选功率与电流均满足要求的型号，则需要先建立型号、供电条件与量值之间的对应，再执行联合筛选。本文据此将领域任务扩展为完整组合查询，使模型提交的程序直接产生包含数值、条件和来源的答案表。

为此，对已有502条来源读法建立类型视图。原始文字、主体、条件和出处完整保留，计算字段增加标量、界限、范围、选项、配对值、近似值与容差等类型。严格界限保留比较符，范围保留端点及其包含关系，脉宽与PRF以对应行保存，典型值和最大值通过量值角色区分。多个型号共享的读法仍由一条来源记录表达。该视图连接资料阅读与程序操作：查询读取原文的结构化含义，结果同时携带原始记录，支持逐项核查。

任务组织包含简单属性查询、限定条件选择、多参数筛选、范围适用性判断、空结果和资料未确立的联合配置。答案中，空表保存实际检索范围及筛选条件；联合配置结论保存各参数的独立范围和来源限制。模型提交的是执行所得对象，最终评分据此核对完整内容及其依据。

\subsection{训练、开发与新来源评价划分}

训练材料采用96个虚构来源实例，每个实例包含三个主体、18条读法和六类问题，共形成576条成功参考轨迹。另以16个独立虚构实例构造96道开发题，并加入八来源资料的12道真实组合开发题。合成数值由固定随机过程生成，独立检查器从生成时保存的事实计算答案，再与程序结果比较。这样，训练内容覆盖操作组合与类型处理，真实资料用于检验专业参数的使用。

新增评价材料选取船用雷达、交通雷达和雷达水位传感器资料，共四个来源家族、八个完整规格页、156条读法和24道问题，每个家族六题，见表\ref{tab:offline-sources}。具体来源为FURUNO DRS4W、OTT RLS、smartmicro UMRR-11 Type45和Campbell CS475A\cite{source-v3-furuno-drs4w,source-v3-ott-rls,source-v3-smartmicro-umrr11,source-v3-campbell-cs475a}。这些资料共同呈现主体、工作状态、量值边界与规格表对应关系，同时提供不同应用类别的文档结构。

\begin{table}[htbp]
\centering\small
\caption{新增真实来源评价材料}
\label{tab:offline-sources}
\OfflineSourcesTable
\end{table}

资料选取后先固定完整页面范围，再整理读法和编写问题。另一位AI审阅者独立阅读原件、作答并封存结果，随后逐题与参考比对。24题的答案内容与定位均通过复核，累计检查\OfflineReviewedCitations{}处引用片段；两项关于评分范围和表头引用的建议以新版本落实。参考类型继续登记为来源约束的AI参考。

\begin{table}[htbp]
\centering\small
\caption{组合查询数据的来源隔离划分}
\label{tab:offline-split}
\OfflineSplitTable
\end{table}

表\ref{tab:offline-split}以完整来源实例或真实来源家族为划分单位，来源组数统计各部分知识目录的覆盖范围。开发目录的24组由16组合成实例和八组真实资料组成；790条读法由288条合成读法和502条真实读法组成。12道真实开发题使用其中七个家族，完整目录同时保留第八个家族的资料。训练、开发、评价分别提供独立的完整知识目录，来源组、记录标识和主体均无交集。目录保留同组全部事实及竞争参数。题面附带输出表名、列名和对象类型，用于统一提交格式；数值、答案行数和目标记录由实际查询确定。全部708条参考程序已在隔离后的目录中重新执行，答案内容及来源与划分前一致。
